\section{Problem Setup}
\label{sec:problem-setup}

We study low-latency \task on \dataset \cite{lei2020tvr} in two related evaluation settings. The lightweight subtitle-window baselines in this repo use a clip-conditioned setup: each query is paired with a video clip $\video$ identified by \texttt{vid\_name}, together with aligned subtitle text and a gold temporal span $(t^{*}_{st}, t^{*}_{ed})$. In that setting, the task is to predict the start and end timestamps $(\hat{t}_{st}, \hat{t}_{ed})$ of the matching moment within the given clip. For \xml, we also use the official TVR corpus-level evaluation, where the system retrieves relevant moments from a video corpus rather than only from a known clip. In the final analysis, the stronger XML results are obtained in the VCMR setting rather than in SVMR.

\paragraph{Temporal search space.}
We represent each clip as an ordered sequence of aligned temporal units and associated subtitles. From the subtitle stream, we construct overlapping candidate windows $\window_k = [a_k, b_k]$ with corresponding text. These windows serve as the cheap search units for text-first retrieval. A text-only method may rank these windows and return the best one directly. In the XML-based pipeline, the ranked windows instead act as candidate generators that prune the temporal context before cross-modal localization is applied. The final system output remains a temporal span in the original timeline rather than only a window identifier.

\paragraph{Fixed-feature setting.}
Our work operates in the feature regime provided by \dataset and the official XML implementation. Video and subtitle features required by XML are assumed to be prepared offline, and subtitle-window construction is also treated as offline preprocessing. As a result, the online cost of the system is dominated by query-time window scoring, candidate selection, and any subsequent multimodal localization. This setup is important because the goal of the project is not to reduce feature-extraction cost from raw video, but to reduce the inference cost of temporal search and localization under a fixed set of precomputed features.

\paragraph{Objective.}
The practical objective is to preserve strong multimodal localization quality while reducing online latency. Text retrieval methods such as \bm, dense sentence-transformer retrieval, and their fusion are evaluated primarily as inexpensive subtitle-search baselines and pruning mechanisms. The main question is whether transcript-based pruning can remove enough irrelevant temporal context to make \xml substantially cheaper, while keeping most of its corpus-level retrieval and localization quality.

\paragraph{Evaluation target.}
We report localization quality using intersection-over-union (IoU) between the predicted span and the gold span. For the subtitle-window baselines, the repo evaluates recall at top-$1$ and top-$3$ under IoU thresholds of $0.3$ and $0.5$, along with mean IoU and absolute start and end errors for the top prediction. For official XML runs, we use the TVR standalone evaluation for VCMR, which reports recall at multiple cutoffs under temporal IoU thresholds. Runtime is reported in milliseconds per query, including average and tail latency statistics such as p50 and p95, and when available, stage-specific times for \bm, dense retrieval, and fusion. Although the original TVR paper discusses modality-type labels, we do not use query-type-specific validation breakdowns in this report.
